\section{Случайные величины и векторы}
\label{sec:random-variables-vectors}

\subsection{Случайная величина и её распределение}

Пусть задано вероятностное пространство $(\Omega,\mathcal F,P)$. \emph{Случайной
величиной} называется измеримое отображение
\[
    X:\Omega\to\mathbb R.
\]
Измеримость означает, что для любого борелевского множества $B\subset\mathbb R$
событие $\{\omega:X(\omega)\in B\}$ принадлежит $\mathcal F$. Поэтому случайная
величина переносит вероятность с пространства элементарных исходов на числовую
ось и тем самым задаёт своё распределение
\[
    P_X(B)=P\{X\in B\}.
\]

Универсальным способом задания распределения является \emph{функция
распределения}
\[
    F_X(x)=P\{X<x\}.
\]
Она не убывает, принимает значения от $0$ до $1$ и удовлетворяет
\[
    \lim_{x\to-\infty}F_X(x)=0,\qquad
    \lim_{x\to+\infty}F_X(x)=1.
\]
В пособии СПбПУ используется именно соглашение $F_X(x)=P\{X<x\}$.
При таком определении функция распределения непрерывна слева (в распространённом
альтернативном соглашении $P\{X\leq x\}$ она была бы непрерывна справа);
вероятность скачка в точке выражается как
\[
    P\{X=x\}=F_X(x+0)-F_X(x).
\]

\subsection{Дискретные и непрерывные случайные величины}

Для дискретной случайной величины достаточно указать возможные значения $x_i$
и их вероятности
\[
    p_i=P\{X=x_i\},\qquad p_i\geq0,\qquad \sum_i p_i=1.
\]
Тогда для любого множества значений вероятность находится суммированием
соответствующих $p_i$.

Если распределение абсолютно непрерывно, существует \emph{плотность}
$f_X(x)\geq0$, для которой
\[
    F_X(x)=\int_{-\infty}^{x} f_X(t)\,dt,
    \qquad
    \int_{-\infty}^{\infty}f_X(x)\,dx=1.
\]
Следовательно,
\[
    P\{a\leq X<b\}=\int_a^b f_X(x)\,dx.
\]
Для непрерывного распределения вероятность отдельного значения равна нулю.

\subsection{Математическое ожидание, дисперсия и моменты}

Математическое ожидание характеризует среднее положение распределения. Для
дискретной и непрерывной случайной величины соответственно
\[
    M X=\sum_i x_i p_i,
    \qquad
    M X=\int_{-\infty}^{\infty}x f_X(x)\,dx,
\]
если соответствующие сумма или интеграл существуют.

Дисперсия характеризует разброс относительно среднего:
\[
    D X=M\bigl[(X-MX)^2\bigr]
       =M(X^2)-(MX)^2,
    \qquad
    \sigma_X=\sqrt{DX}.
\]
Начальные и центральные моменты порядка $k$ определяются как
\[
    m_k=M(X^k),\qquad
    \mu_k=M\bigl[(X-MX)^k\bigr].
\]
В частности, $m_1=MX$, а $\mu_2=DX$.

Если $Y=g(X)$, то при известном законе распределения $X$ числовые
характеристики функции случайного аргумента можно вычислять без предварительного
нахождения закона $Y$:
\[
    M g(X)=\int_{-\infty}^{\infty}g(x)f_X(x)\,dx
\]
в непрерывном случае и аналогичной суммой в дискретном.

\subsection{Случайный вектор и совместное распределение}

\emph{Случайным вектором} называется набор случайных величин
\[
    \mathbf X=(X_1,\ldots,X_n)^T.
\]
Его распределение является вероятностным распределением в $\mathbb R^n$.
Совместная функция распределения может быть записана как
\[
    F_{\mathbf X}(x_1,\ldots,x_n)
    =P\{X_1<x_1,\ldots,X_n<x_n\}.
\]
Для абсолютно непрерывного вектора вводится совместная плотность
$f_{\mathbf X}(x_1,\ldots,x_n)$, причём вероятность попадания случайной точки
в область $G\subset\mathbb R^n$ равна
\[
    P\{\mathbf X\in G\}
    =\int_G f_{\mathbf X}(\mathbf x)\,d\mathbf x.
\]

Одномерные, или маргинальные, распределения получают интегрированием совместной
плотности по остальным переменным. Например, для пары $(X,Y)$
\[
    f_X(x)=\int_{-\infty}^{\infty}f_{X,Y}(x,y)\,dy.
\]
Случайные величины $X$ и $Y$ независимы тогда и только тогда, когда их
совместный закон факторизуется; для плотностей это означает равенство почти всюду
\[
    f_{X,Y}(x,y)=f_X(x)f_Y(y).
\]

\subsection{Условное распределение и условное математическое ожидание}

Если $f_X(x)>0$, условная плотность $Y$ при фиксированном $X=x$ имеет вид
\[
    f_{Y|X}(y|x)=\frac{f_{X,Y}(x,y)}{f_X(x)}.
\]
Условное математическое ожидание
\[
    M(Y\mid X=x)=\int_{-\infty}^{\infty}y f_{Y|X}(y|x)\,dy
\]
характеризует среднее значение $Y$ при известном значении $X$. В корреляционной
теории эта зависимость называется регрессией $Y$ на $X$.

Для случайного вектора удобно вводить вектор математических ожиданий
\[
    \mathbf m=M\mathbf X=(MX_1,\ldots,MX_n)^T.
\]
Характеристики взаимной зависимости его компонент --- ковариации и
корреляционная матрица --- являются предметом следующего вопроса.

\subsection{Предельные теоремы}

Две базовые предельные теоремы связывают поведение большой совокупности
независимых испытаний с детерминированным средним и нормальным законом.

\textbf{Слабый закон больших чисел (в простейшей форме).}
Пусть $X_1,X_2,\ldots$ независимы и одинаково распределены,
\[
    MX_k=\mu,\qquad DX_k=\sigma^2<\infty.
\]
Тогда выборочное среднее
\[
    \overline X_n=\frac1n\sum_{k=1}^nX_k
\]
сходится к $\mu$ по вероятности:
\[
    \overline X_n\xrightarrow{P}\mu.
\]
Действительно, из неравенства Чебышёва следует
\[
    P\{|\overline X_n-\mu|\geq\varepsilon\}
    \leq \frac{\sigma^2}{n\varepsilon^2}
    \longrightarrow0.
\]

\textbf{Центральная предельная теорема Линдеберга--Леви.}
Если дополнительно $0<\sigma^2<\infty$, то
\[
    \frac{\sum_{k=1}^nX_k-n\mu}{\sigma\sqrt n}
    \xrightarrow{d}N(0,1),
\]
то есть для любой точки непрерывности функции распределения стандартного
нормального закона
\[
    P\left\{
      \frac{\sum_{k=1}^nX_k-n\mu}{\sigma\sqrt n}<x
    \right\}
    \longrightarrow \Phi(x).
\]
Закон больших чисел описывает стабилизацию среднего, а центральная предельная
теорема --- форму случайных отклонений суммы от её среднего значения.

\subsection{Что важно сказать на экзамене}

Нужно начать с определения случайной величины как измеримого отображения и
объяснить, что её закон задаётся функцией распределения. Затем различить
дискретный закон и непрерывный закон с плотностью, дать формулы математического
ожидания и дисперсии. Для случайного вектора важно определить совместное
распределение, показать, как из него получают маргинальные и условные
распределения, и сформулировать критерий независимости через факторизацию
совместной плотности. Условное математическое ожидание удобно интерпретировать
как среднее значение одной величины при известном значении другой. В качестве
основных предельных результатов достаточно уверенно сформулировать слабый закон
больших чисел и центральную предельную теорему для независимых одинаково
распределённых величин.

\newpage
